\section{Related Work and Motivation}
\label{sec:related}

\subsection{Bone Suppression and CT-Derived Supervision}
Early learning-based rib suppression used massive-training artificial neural networks \cite{suzuki2006mtann}. Later work used
convolutional networks, for example a cascade of multi-scale networks in the gradient domain \cite{yang2017cascade}. Recent methods train
diffusion models on paired radiographs and DES soft-tissue images
\cite{chen2024bsdiff,sun2026bsldm,sun2025gllcm,sun2026debonedit}; the SZCH-X-Rays data set was compiled for BS-LDM \cite{sun2026bsldm}
and made public with DeBoneDiT \cite{sun2026debonedit}. Paired DES images are scarce and come from a few scanners. CT removes this
bottleneck: Angaitkar et al. \cite{angaitkar2026anatomy} build digitally reconstructed radiographs from chest CT, project bone, non-lung
soft tissue and lung separately, and train models that predict one component. Their projections are released and are built from CT-RATE
volumes \cite{hamamci2026ctrate}. We train on these projections and use the released weights as the main comparator. Our models differ in
encoder and resolution (ConvNeXt-V2-T \cite{woo2023convnextv2} at 512~px against SwinV2-B \cite{liu2022swinv2} at 1024~px) and in
sample budget, and they add the twin regulariser, the ensemble and the audit.

\subsection{Lesion Fidelity and Counterfactual Auditing}
Learned image-to-image models can remove or invent content: distribution matching in translation networks can drop or hallucinate lesions
\cite{cohen2018distribution}, deep reconstruction can be unstable to small structures \cite{antun2020instabilities}, and hallucinations
of learned reconstruction have been defined and measured \cite{bhadra2021hallucinations,tivnan2024hallucination}. Inserting known
features or signals into images is an established way to test learned MRI reconstruction \cite{cheng2020addressing} and the effect of
learned denoising on signal detection \cite{li2021assessing}. For bone suppression, simulated nodules were inserted to measure contrast
after a hand-built clavicle suppression method \cite{hogeweg2013suppression}, and a phantom with a simulated nodule at many positions was
used to measure contrast after bone suppression \cite{takagi2020dose}. These studies measure lesion contrast. Contrast alone cannot
separate a suppressor that keeps a lesion from one that suppresses weakly, and a score computed over the whole image is dominated by
normal anatomy. Our audit therefore measures the fraction of a known additive lesion signal that a learned suppressor moves into its bone
image, on real radiographs, at sites shared by every model, and reports it next to suppression strength.

\subsection{Suppressed Images as Inputs, Explanations and Uncertainty}
Suppressed images have been used to classify tuberculosis-consistent findings \cite{rajaraman2021bone}, and readers and computer-aided
detection have been compared with and without them \cite{freedman2011lung,schalekamp2014cad}. A suppressed image can be given to a
classifier alone, stacked with the radiograph as channels, or as a separate input stream; the last two add computation, which a fair
comparison has to match with a control of equal capacity. For multi-label classification we use a class-query decoder over spatial tokens
\cite{liu2021query2label,ridnik2023mldecoder}, so that the attention of each class on each input stream can be read out, and attribution
maps from Grad-CAM++ \cite{chattopadhay2018gradcampp} and Integrated Gradients \cite{sundararajan2017axiomatic} scored with the pointing
game \cite{zhang2018topdown}. For uncertainty we use deep ensembles \cite{lakshminarayanan2017ensembles}, split-conformal prediction
\cite{lei2018conformal,angelopoulos2023conformal}, conformal risk control of the false-negative rate \cite{angelopoulos2022crc},
class-conditional (Mondrian) sets \cite{vovk2005algorithmic} and selective classification \cite{geifman2017selective,geifman2019bias}.

Table~\ref{tab:literature} compares the studies closest to ours. Three gaps motivate the design: lesion erasure by learned suppressors
has not been measured against an exactly known signal; gains from suppressed inputs have not been compared against a capacity-matched
control; and the uncertainty and localisation of CT-supervised suppression have not been examined together with its downstream benefit.

\begin{table*}[!t]
\centering
\caption{Comparison of studies on bone suppression and suppressed images as model inputs with the present work. Check marks indicate
whether a study reports each element: a counterfactual or simulated-lesion audit (Lesion), a capacity-matched control (Cap.),
uncertainty quantification or calibration (UQ), validation on external data (Ext.) and explanation maps (XAI). \checkmark: reported;
(\checkmark): qualitative only; --: not reported; n/a: not applicable to the study design (no model trained by the authors, no input added
to a model, or no clinical task). Rows marked $^{\dagger}$ were assessed from the abstract because the full text is not openly accessible;
all other rows from the full text.}
\label{tab:literature}
\footnotesize
\setlength{\tabcolsep}{3pt}
\begin{tabularx}{\textwidth}{@{}>{\raggedright\arraybackslash}p{2.6cm}>{\raggedright\arraybackslash}p{3.4cm}%
>{\raggedright\arraybackslash}p{3.3cm}>{\raggedright\arraybackslash}X*{5}{c}@{}}
\toprule
Study & Data & Method & Key result & Lesion & Cap. & UQ & Ext. & XAI\\
\midrule
Suzuki et al.$^{\dagger}$ \cite{suzuki2006mtann} & DES bone images for training; 118 validation and 136 independent test radiographs (14
 institutions) & Multiresolution massive-training ANN (MTANN) & Ribs and clavicles suppressed while nodule and vessel visibility was kept
 & -- & n/a & -- & \checkmark & n/a\\
Yang et al.$^{\dagger}$ \cite{yang2017cascade} & 504 two-exposure DES radiographs (404 training, 100 test) & Cascade of multi-scale CNNs
 in the gradient domain (CamsNet) & Soft-tissue PSNR 38.7\,dB, bone suppression ratio 83.8\,\%; other X-ray machines shown visually & --
 & n/a & -- & (\checkmark) & n/a\\
Chen et al. \cite{chen2024bsdiff} & Paired radiographs and DES soft tissue of 120 patients, one hospital (7:2:1 split) & Conditional
 diffusion model (BS-Diff) & Diffusion-based suppression with a clinical image-quality evaluation & -- & n/a & -- & -- & n/a\\
Sun et al. \cite{sun2026bsldm} & SZCH-X-Rays and JSRT (separate splits); Shenzhen set for a downstream test & Conditional latent
 diffusion (BS-LDM) & High-resolution suppression; clinical and downstream lesion-detection evaluation & -- & n/a & -- & \checkmark
 & n/a\\
Sun et al. \cite{sun2025gllcm} & SZCH-X-Rays (592/74/75) and JSRT (192/24/25) & Global-local latent consistency model (GL-LCM) &
 Suppression in 50 sampling steps, 8.54\,s against 84.62\,s for BS-LDM & -- & n/a & -- & -- & n/a\\
Sun et al. \cite{sun2026debonedit} & SZCH-X-Rays (released here) and JSRT; Asraf set for a downstream test & Depth-driven conditional
 bridge diffusion transformer (DeBoneDiT) & Public paired DES data; downstream classification on an external set & -- & n/a & -- &
 \checkmark & n/a\\
Angaitkar et al. \cite{angaitkar2026anatomy} & \num{21887} CT-RATE projections; TBX11K, Node21, VinDr-CXR & Anatomy-decomposed CT
 projections, SwinV2-B at 1024~px & Detection gains and higher real-nodule contrast; no lesion-insertion audit & -- & n/a & -- &
 \checkmark & --\\
Hogeweg et al.$^{\dagger}$ \cite{hogeweg2013suppression} & 20 simulated radiographs; 253 radiographs; 100 with CT nodules simulated
 near the clavicles & Model-based suppression of elongated structures & Simulated nodules kept more contrast than the clavicles after
 suppression (1.35 vs.\ 0.55) & \checkmark & n/a & -- & n/a & n/a\\
Takagi et al.$^{\dagger}$ \cite{takagi2020dose} & Anthropomorphic phantom, 12\,mm simulated nodule at 66 positions & Commercial
 bone-suppression processing at low tube voltage & Nodule contrast and CNR rose away from the ribs; lower dose possible & \checkmark &
 n/a & -- & n/a & n/a\\
Freedman et al.$^{\dagger}$ \cite{freedman2011lung} & Reader study, 368 patients, 15 radiologists & Rib and clavicle suppression software
 & Reader LROC AUC rose from 0.460 to 0.558 & -- & n/a & n/a & n/a & n/a\\
Schalekamp et al.$^{\dagger}$ \cite{schalekamp2014cad} & Reader study, 300 radiographs from four institutions, 8 readers & CAD with
 bone-suppressed images available & CAD raised reader AUC from 0.812 to 0.841 beyond suppressed images & -- & n/a & n/a & n/a & n/a\\
Rajaraman et al. \cite{rajaraman2021bone} & JSRT and NIH-CC DES (suppression); Shenzhen and Montgomery TB sets & Suppressed image
 replacing the radiograph as classifier input (VGG-16) & TB AUC 0.899 to 0.954 (Shenzhen) with suppressed inputs; cross-institutional
 suppression test & -- & n/a & -- & \checkmark & \checkmark\\
\midrule
This work & Qure.ai projections; SZCH-X-Rays, VinDr-RibCXR, JSRT, Node21, TBX11K, VinBigData, NIH ChestX-ray14; BDCXR-3257 & Twin
 lesion-invariance regulariser, $K=3$ ensemble; suppressed images as extra channels; multi-stream classifier with a capacity-matched
 control & $A=0.051$ vs.\ 0.149 (released) and 0.178 (control); Node21 CPM $+0.028$ with extra channels
 & \checkmark & \checkmark & \checkmark & \checkmark & \checkmark\\
\bottomrule
\end{tabularx}
\end{table*}
